% ==========================================
% 1.3 State of The Art - Software Development
% ==========================================
 
\vspace*{0.5cm}
\subsection{State of the Art in Software Development}\label{sub:sota}
 
Modern software development has undergone a structural transformation
over the past decade, driven by four converging forces: the rise of
AI-assisted tooling, the maturation of DevOps and CI/CD practices,
the widespread adoption of cloud-native architectures, and the
continued evolution of Agile methodologies.
This section surveys the current state of the art in each of these
areas, establishing a baseline against which the practices observed
during the internship are later contrasted in
Section~\ref{sub:contrast}.
 
 
\subsubsection{AI-Assisted Developer Tooling}\label{subsub:sota:ai}
 
The integration of large language models into the software development
workflow has moved from experimental to mainstream.
Tools such as GitHub Copilot, powered by OpenAI's Codex, now act as
context-aware coding assistants, offering real-time code suggestions
and automating repetitive tasks directly inside the developer's
IDE~\cite{softjourn2025aidevops}.
Beyond individual productivity, generative AI has expanded into every
phase of the development lifecycle: AI-driven requirements synthesis
and architecture generation at the planning stage, autonomous
test-case creation during quality assurance, and self-healing anomaly
detection during maintenance~\cite{cio2026future}.
 
The Perforce \textit{2026 State of DevOps Report}, surveying 820
technology professionals worldwide, confirms that AI now permeates
software delivery at scale~\cite{perforce2026devops}.
It frames AI not as a replacement for DevOps discipline, but as an
\textit{amplifier} of existing organisational states: 72\,\% of
leaders in high-maturity DevOps organisations report deeply embedded
AI practices, compared to just 18\,\% in low-maturity
counterparts~\cite{perforce2026devops}.
Despite enthusiasm, governance remains fragile: only 39\,\% of
organisations maintain fully automated audit trails, and trust in
AI output is far from universal~\cite{perforce2026devops}.
A complementary Harness report from March 2026, based on 700
engineering practitioners across five countries, further cautions that
AI coding tools have dramatically accelerated code production, but
the downstream delivery pipeline, encompassing testing, security
scanning, and deployment, has not kept pace~\cite{harness2026devops}.
 
 
\subsubsection{DevOps and Continuous Integration / Continuous
Delivery}\label{subsub:sota:devops}
 
DevOps has become the default operating model for software delivery
organisations. In 2026 the field is characterised by three prominent
evolutions:
 
\begin{itemize}
    \item \textbf{AI-augmented operations (AIOps).}
    AI is embedded in monitoring, incident management, and deployment
    pipelines. Predictive analytics identify failure patterns before
    incidents occur, while self-healing systems detect anomalies and
    apply automated remediations without human
    intervention~\cite{devopsdotcom2025}. Gartner projects that
    73\,\% of enterprises will be implementing or planning AIOps
    adoption by end of 2026~\cite{fullstack2026aidevops}.
 
    \item \textbf{DevSecOps and shift-left security.}
    Security is integrated throughout the CI/CD pipeline rather than
    validated at the end. Automated vulnerability scanning and static
    application security testing (SAST) run on every commit, making
    security a shared responsibility across development, operations,
    and security teams~\cite{graphite2025}. As of 2026, 97\,\% of
    organisations report using or planning to use AI in the software
    delivery lifecycle, making AI governance a critical component of
    any DevSecOps strategy~\cite{devsecops2026}.
 
    \item \textbf{Platform engineering and GitOps.}
    GitOps treats the Git repository as the single source of truth
    for infrastructure state. Gartner predicts that 80\,\% of
    software development companies will adopt Internal Developer
    Platforms by 2026 to unify tooling, automate governance, and
    accelerate delivery at scale~\cite{fullstack2026aidevops}.
    Yet the Harness 2026 report finds that 73\,\% of engineering
    leaders say hardly any development teams have standardised
    templates or golden paths for services and pipelines, indicating
    that adoption and maturity remain
    uneven~\cite{harness2026devops}.
\end{itemize}
 
The Perforce 2026 report introduces a sharp operational divide: mature
organisations are 36\,\% more likely to automate the majority of
deployments from commit to production, and 66\,\% more likely to
respond very effectively to production incidents using automated
rollbacks~\cite{perforce2026devops}. This underscores that
the value of DevOps investments is increasingly contingent on
organisational maturity rather than tool adoption alone.
 
 
\subsubsection{Cloud-Native Architectures and
Microservices}\label{sec:sota:cloud}
 
Cloud-native development has transitioned from competitive advantage
to operational baseline.
The 2025 CNCF Annual Cloud Native Survey, published in January 2026
and covering 628 IT professionals across industries, reports that
98\,\% of organisations now use some form of cloud-native technology,
and 82\,\% of container users run Kubernetes in production, up from
66\,\% in 2023~\cite{cncf2026survey}.
Kubernetes now functions as the common denominator for AI
workloads as well: 66\,\% of organisations hosting generative AI
models use Kubernetes to manage some or all of their inference
workloads~\cite{cncf2026survey}.
 
The architectural centrepiece of this paradigm remains the
decomposition of monolithic applications into small, loosely coupled
\textit{microservices}, each independently deployable and scalable.
Looking ahead, 84\,\% of enterprises expect to build at least half
of their new applications on Kubernetes within five years, signalling
a clear transition towards container-first development
models~\cite{portworx2026}.
A complete cloud-native stack extends beyond Kubernetes to include
CI/CD pipelines, service meshes (e.g.\ Istio or Linkerd),
distributed tracing, and platform engineering practices that abstract
infrastructure complexity from application
developers~\cite{ini8labs2025}.
 
Resilience patterns, such as Circuit Breaker, Retry with exponential
backoff, and Health Checks, are now standard design primitives used
to prevent cascading failures in distributed
systems~\cite{madrigan2025}.
 
 
\subsubsection{Agile Methodologies and Team
Practices}\label{subsub:sota:agile}
 
Agile practices remain the dominant framework for software project
management, though their application has matured considerably.
The Perforce 2026 report identifies DevOps maturity as a
prerequisite for AI scalability rather than a practice made obsolete
by it: organisations with disciplined engineering, automation, and
strong collaboration are the ones scaling AI successfully and
turning innovation into measurable business
outcomes~\cite{perforce2026devops}.
The report frames this as a \textit{shift-up} in roles: 87\,\% of
respondents believe AI will enable engineers to focus less on
scripting and more on system design and directing
outcomes~\cite{perforce2026devops}.
 
Contemporary high-performing engineering organisations combine Agile
sprints with platform engineering initiatives providing developers
with self-service golden paths, internal developer portals, and
standardised toolchains.
However, the Harness 2026 report reveals a significant gap between
aspiration and reality: the majority of surveyed developers and their
managers acknowledge that their current ways of working will not be
sustainable over the long term, with the figure rising to 81\,\% for
very frequent users of AI coding tools~\cite{harness2026devops}.
This tension, between acceleration at the code level and inadequate
automation downstream, represents the central challenge for Agile
organisations in 2026.
 
\bigskip
\noindent
The practices described above represent the current industry standard
for well-resourced software organisations in 2025--2026. The
following section contrasts these benchmarks with the tooling,
processes, and architectural choices encountered during the
internship, highlighting both areas of alignment and notable gaps.
 
 
% ============================================================
%  Comparative Internship Practices - Software Development
% ============================================================
\vspace*{0.5cm}
\subsection{Comparative Internship Practices}\label{sub:contrast}
 
The internship project, while built by a small team under no
client-constraints, intersects with many of those benchmarks in
meaningful ways and diverges from others in ways that are equally
instructive.
The following subsections examine each dimension in turn.
 
 
\subsubsection{Architecture and Design Patterns vs. the Cloud-Native
Paradigm}\label{subsub:contrast:arch}
 
The state of the art favours decomposing applications into
independently deployable microservices orchestrated by Kubernetes,
with service meshes handling inter-service communication and
resilience~\cite{ini8labs2025,madrigan2025}.
The internship project took a different, yet deliberate, approach:
a \textit{layered N-Tier architecture} partitioned into discrete
projects (\texttt{API}, \texttt{Services}, \texttt{Repositories},
and shared utility layers) running as a single deployable backend
unit.
 
This is a conscious alignment with the \textit{modular monolith}
pattern, which many practitioners recognise as a pragmatic
starting point for teams that have not yet reached the scale that
justifies the operational overhead of a full microservices
deployment~\cite{ini8labs2025}.
The strict separation of concerns enforced by the Repository Pattern
and Dependency Injection mirrors the loose-coupling principle at the code level, that microservices architecture seeks to achieve at the infrastructure level.
The use of Factory and Strategy/Handler patterns to replace branching
logic is likewise consistent with maintainability principles that are
independent of deployment topology.
 
In terms of container adoption, the project is well-aligned with
industry practice. The backend is fully containerised: the CI/CD
pipeline builds a Docker image on every push, tags it with the
commit hash for traceability, pushes it to \textbf{AWS Elastic
Container Registry (ECR)}, and deploys it via \textbf{AWS Elastic
Container Service (ECS)}. This satisfies the core cloud-native
principle of reproducible, immutable deployments~\cite{cncf2026survey}.
 
The only distinction relative to the state of the art is the choice
of \textbf{AWS ECS} over \textbf{Kubernetes}. While 82\,\% of
container users now run Kubernetes in production~\cite{cncf2026survey},
ECS is a fully managed container orchestration service that handles
scheduling, scaling, and health management without the operational
overhead of self-managing a Kubernetes cluster. For a project at this
scale, ECS represents a pragmatic and legitimate trade-off: it
delivers the essential benefits of container orchestration while
reducing infrastructure complexity, a pattern explicitly encouraged
by the CNCF, which notes that 79\,\% of Kubernetes users themselves
opt for managed services over self-managed
clusters~\cite{cncf2026survey}.
 
 
\subsubsection{Resilience and Observability vs. Industry
Standards}\label{subsub:contrast:resilience}
 
The state of the art calls for resilience patterns to be implemented
as standard design primitives in distributed systems, alongside
AIOps-driven observability for anomaly detection and
self-healing~\cite{madrigan2025,devopsdotcom2025}.
 
The internship project demonstrates a solid alignment with the
resilience dimension: outbound HTTP calls are wrapped in automatic
retry policies with exponential backoff using the \textbf{Polly}
library, directly implementing the pattern described in the
literature. Error observability is addressed through a centralised
\textbf{Sentry} integration, which provides real-time crash
reporting and structured error tracking, a practice consistent with
the industry emphasis on proactive incident
management~\cite{devopsdotcom2025}.
 
The gap here lies in the scope of observability.
Industry-leading deployments combine distributed tracing, structured
logging, and AI-assisted anomaly detection across the full
application stack. The Perforce 2026 report identifies the absence of
automated audit trails as a critical governance vulnerability, with
only 39\,\% of organisations currently maintaining
them~\cite{perforce2026devops}. The internship implementation covers
the error layer effectively but does not extend to distributed
tracing or performance profiling, areas that would become critical as
the system scales.
 
 
\subsubsection{Data Strategy vs. Cloud-Native Data
Practices}\label{subsub:contrast:data}
 
Modern cloud-native architectures advocate for polyglot persistence,
choosing the right data store for each workload, as well as managed
cloud-native database services that abstract operational
complexity~\cite{techzine2025}.
The internship project explicitly adopts a \textbf{Polyglot
Persistence} strategy, combining \textbf{PostgreSQL} for structured
relational data and \textbf{MongoDB} for semi-structured document
data, which is a direct application of this principle.
 
The choice of \textbf{Dapper} as a lightweight data-access layer
over PostgreSQL is a deliberate trade-off: it forgoes the abstraction
of a full ORM in favour of hand-optimised SQL queries, which aligns
with the industry guidance to maintain performance efficiency in
data-intensive paths.
 
 
\subsubsection{Frontend Architecture vs. Current Mobile Development
Practices}\label{sec:contrast:frontend}
 
Cross-platform mobile development using \textbf{React Native} and the
\textbf{Expo} ecosystem is well-established in the industry, and the
architectural choices made in the internship project are broadly
consistent with current best practices.
The feature-based folder structure, the use of TypeScript for strict
typing, and the isolation of data-fetching into dedicated service
functions using \textbf{Axios} all reflect the separation-of-concerns
principle that underpins modern frontend architecture.
 
The decision to use React's native \textbf{Context API} rather than a
heavyweight global state manager represents a pragmatic
fit-for-purpose choice, appropriate for the application's scope and
consistent with the current community trend of preferring lighter,
composable state solutions over monolithic state libraries.
 
A notable gap is the \textbf{offline-first} capability.
Offline-first is increasingly considered a quality baseline for mobile
applications in the industry, particularly those involving location
and real-time data. The feature was implemented but subsequently
broken by upstream library updates from React Native and Expo, and
moved to the backlog. This outcome illustrates a real and widely
reported challenge in the React Native ecosystem: the pace of
framework updates can introduce unplanned regressions, a maintenance
burden that enterprise teams typically address through stricter
dependency pinning and automated regression testing in CI/CD
pipelines~\cite{devopsdotcom2025}.
 
 
\subsubsection{CI/CD Maturity and DevOps
Tooling}\label{subsub:contrast:aidevops}
 
The state of the art positions automated CI/CD pipelines as
table-stakes for professional software development teams in
2026~\cite{perforce2026devops,harness2026devops}.
The internship project presents a mixed but largely encouraging
picture in this dimension, with a meaningful distinction between the
backend and frontend repositories.
 
The backend pipeline is a strong alignment with industry practice.
Implemented in Bitbucket Pipelines, it follows a multi-stage
structure: feature branches trigger automated tests and a Docker
image build on every push; merges to \texttt{master} additionally
promote the image through a full environment pipeline, automatically
deploying to a \textbf{staging} environment via AWS ECS, and then
offering a \textbf{manual trigger} for the production deployment.
This constitutes a complete \textit{CI/CD pipeline with environment
promotion and a gated production release}, which is precisely the
pattern recommended for production systems~\cite{perforce2026devops}.
The commit-tagged Docker images pushed to AWS ECR further provide
full deployment traceability, enabling rollback to any prior
revision. The manual gate on production is a deliberate
human-in-the-loop safety check rather than an absence of automation,
and is standard practice even in high-maturity
organisations~\cite{perforce2026devops}.
 
The frontend repository presents a more significant gap: it lacked
automated tests entirely, meaning the CI safety net that existed on
the backend had no equivalent on the client side. The Harness 2026
report identifies the absence of standardised pipelines and
automated quality gates as one of the primary reasons AI-accelerated
code production leads to downstream instability~\cite{harness2026devops}.
While the frontend was not generating code at AI-assisted velocity,
the underlying risk is the same: without automated regression
coverage, regressions introduced by dependency updates (such as the
offline-first breakage described in Section~\ref{sec:contrast:frontend})
are harder to detect systematically and earlier in the development
cycle.
 
\newpage
\subsubsection{Summary}\label{subsub:contrast:summary}
 
\begin{table}[ht]
\centering
\small
\renewcommand{\arraystretch}{1.4}
\begin{tabular}{@{} >{\centering\arraybackslash}m{3.2cm}
                    >{\centering\arraybackslash}m{5.8cm}
                    >{\centering\arraybackslash}m{5.8cm} @{}}
\toprule
\textbf{Dimension} & \textbf{Alignment} & \textbf{Gap} \\
\midrule
 
Architecture &
  Layered N-Tier; loose coupling via DI and Repository Pattern;
  containerised with Docker; deployed via AWS ECS &
  ECS used instead of Kubernetes; managed service trade-off
  appropriate for project scale \\
 
Resilience \& Observability &
  Retry / exponential backoff (Polly); centralised error
  tracking (Sentry) &
  No distributed tracing or AI-driven anomaly detection \\
 
Data Strategy &
  Polyglot persistence (PostgreSQL + MongoDB); optimised
  queries via Dapper &
  No managed cloud-native database services \\
 
Frontend &
  React Native + Expo; TypeScript; feature-based structure;
  Context API &
  Offline-first capability regressed and not restored \\
 
CI/CD Maturity &
  Backend: full CI/CD pipeline with environment promotion
  (staging auto-deploy, production manual gate), commit-tagged
  images on AWS ECR, deployed to AWS ECS &
  Frontend: no automated tests or CI pipeline \\
 
\bottomrule
\end{tabular}
\caption{Alignment of internship practices with the state of the art.}
\label{tab:contrast}
\end{table}
 
\noindent
Overall, the internship project demonstrates a strong grasp of
foundational software engineering principles (clean architecture,
resilience, and polyglot data management) that are prerequisites for
any cloud-native or AI-augmented evolution.
The gaps identified are primarily operational and infrastructural
rather than conceptual, and represent natural next steps as the
project and team mature.



\subsection{Things used}
\subsubsection{Backend Architecture and Project Structure}
For the server side, the project uses a classic layered approach (N-Tier architecture) to keep the code clean and separated. Instead of putting everything into one giant folder, the backend is organized into multiple specialized projects, each handling a specific job:
\begin{itemize}
    \item \textbf{\texttt{API} (Presentation):} This is the entry point of the application. It handles the RESTful HTTP endpoints (Controllers), custom middlewares, and general configuration settings.
    \item \textbf{\texttt{Services} (Business Logic):} This layer contains the core rules of the app, data processing logic, and custom data validation classes.
    \item \textbf{\texttt{Repositories} (Data Access):} This layer is solely responsible for talking to the databases, keeping database queries isolated from the business logic.
    \item \textbf{\texttt{Models}, \texttt{Shared}, and \texttt{Tests}:} These tiers store the Data Transfer Objects (DTOs), cross-project utility constants, and automated unit tests to ensure code reliability.
\end{itemize}

\subsubsection{Software Design Patterns Used}
To make the codebase maintainable and easy to expand during development, several industry-standard design patterns were used:
\begin{itemize}
    \item \textbf{Repository Pattern:} Data access is hidden behind interfaces. This means the services don't care how data is saved or retrieved, making it easy to change database mechanics later without breaking the core application.
    \item \textbf{Dependency Injection (DI):} Objects and services are injected through constructors rather than instantiated manually using the \texttt{new} keyword. Using varying lifecycles (Singleton, Scoped, and Transient), the framework manages memory efficiently and keeps components decoupled.
    \item \textbf{Factory and Strategy/Handler Patterns:} To handle different actions dynamically, such as triggering different logic for an \texttt{Achievement} versus a \texttt{Walk} notification, the app combines a Factory (\texttt{NotificationFactory}) with specialized handlers. This replaces hard-to-read, hard-to-maintain \texttt{switch/if} blocks.
    \item \textbf{DTOs and Mappers:} The app uses Data Transfer Objects to define exactly what data goes over the network, ensuring internal domain database structures are never exposed directly to the frontend. Custom mappers are used to easily convert data back and forth between these formats.
    \item \textbf{Async/Await Pattern:} Every layer consistently uses asynchronous programming using \texttt{Task<T>}. This prevents the server from blocking threads while waiting for slow operations like database queries or network responses, massively increasing how many users the API can handle at once.
\end{itemize}

\subsubsection{Polyglot Persistence (Databases)}
This project uses a \textbf{Polyglot Persistence} strategy, meaning it mixes different database types depending on what data needs to be stored, offering the best of both worlds:
\begin{itemize}
    \item \textbf{PostgreSQL} is used as the relational database for structured, transactional data. To interact with it quickly, the project uses \textbf{Dapper}, a lightweight library that bridges the gap between relational databases and object oriented code, allowing for writing fast and optimized native SQL queries.
    \item \textbf{MongoDB} is used as a NoSQL document database to handle more dynamic or semi-structured data fluidly, communicating through the native \texttt{MongoDB.Driver}.
\end{itemize}

\subsubsection{Resiliency, Monitoring, and Error Handling}
When an API talks to external cloud services it has to be prepared for network hiccups or rate limits ($429$ errors). The backend implements the \textbf{Resilience Pattern} using a library called \textbf{Polly}. This wraps outbound HTTP requests in automatic retry policies with exponential backoff, making sure transient failures don't crash the app for the user.

Additionally, the project uses a custom \textbf{Middleware} to catch any unhandled crashes globally. It automatically logs these errors and forwards them to a centralized tracking hub called \textbf{Sentry} so developers can debug them in real-time.

\subsubsection{Frontend Mobile Architecture (React Native \& Expo)}
On the frontend, the client layer is built as a cross-platform mobile app using \textbf{React Native} on top of the \texttt{Expo} ecosystem. The architecture follows a clean, feature-based and functional layout. 

Instead of pulling in a massive, heavy global state manager, the application manages global data using React’s native \textbf{Context API}. State is split into four distinct domains:
\begin{enumerate}
    \item \textbf{\texttt{UserContext}:} Handles user authentication and profile sessions.
    \item \textbf{\texttt{LocationContext}:} Streams and manages real-time GPS coordinates.
    \item \textbf{\texttt{ConnectionContext}:} Monitors current internet and network status.
    \item \textbf{\texttt{NotificationContext}:} Manages push notification tokens.
\end{enumerate}

The frontend's data layer matches the backend's clean separation of concerns. Screen components never fetch data directly; instead, they call isolated service functions built using \textbf{Axios} that return strictly typed TypeScript definitions. Finally, to ensure a smooth user experience, the client implements an \textbf{Offline-First approach}, combining real-time background location tracking with local storage to cache data when the user loses internet connection, but, due to multiple react and expo updates, this feature (which was previously implemented) no longer worked and was sent to the backlog of tasks.


% ============================================================
%  Comparative Internship Practices - Market (to be integrated into Description of Work
% ============================================================
\vspace*{0.5cm}
\subsection{Comparative Internship Practices}\label{sub:contrast}
 
The internship project, while built by a small team under no
client-constraints, intersects with many of those benchmarks in
meaningful ways and diverges from others in ways that are equally
instructive.
The following subsections examine each dimension in turn.

 
\subsubsection{Location-Based Features and Offline-First
Capability}\label{subsub:contrast:location}
 
The internship application is directly positioned in the
location-based outdoor exploration space, using real-time GPS
tracking to trigger trail and POI interactions as users move
through physical space. This aligns with the industry direction
described in Section~\ref{subsub:sota:location}, where physical
presence detection is the core interaction primitive.
 
The most significant gap in this dimension is the \textbf{offline-first}
capability.
As noted in Section~\ref{subsub:sota:location}, offline-first
architecture is considered a baseline expectation for outdoor
trail applications~\cite{markwide2026hiking}, and the
internship project did implement this feature.
However, it was subsequently broken by upstream library updates
from React Native and Expo and moved to the backlog.
This outcome illustrates a real and widely reported challenge in
the React Native ecosystem: the pace of framework updates can
introduce unplanned regressions, a maintenance burden that
enterprise teams typically address through stricter dependency
pinning and automated regression testing in CI/CD
pipelines~\cite{markwide2026hiking}.
 
 
\subsubsection{Gamification Design}\label{subsub:contrast:gamification}
 
The internship application implements a core gamification loop
centred on points awarded for completing trivia-style activities at
POIs along trails, with an achievement system tracking user
milestones.
This is a direct application of the hybrid reward architecture
described in Section~\ref{subsub:sota:gamification}: the
activity-based point system provides intrinsic game progression,
while the physical trail context grounds the rewards in real-world
exploration.
 
The design is well-aligned with the academic findings on
gamification effectiveness~\cite{ama2025gamification}: rewards
are tied to completing a value-added activity (a trivia question
about a real POI) rather than arbitrary engagement actions,
which mitigates the risk of gamification flow decoupling from the
application's core purpose.
 
A gap relative to the state of the art is the absence of
\textit{adaptive personalisation}.
Leading gamification implementations in 2026 use AI-driven
behavioural data to personalise challenge difficulty, reward
cadence, and content selection~\cite{amplifai2026gamification}.
The current system presents uniform content and reward structures
regardless of user history, which represents a natural area for
future development as the platform matures.
 
 
\subsubsection{LLM API Integration for Trivia
Generation}\label{subsub:contrast:llm}
 
The internship application integrates a cloud-based LLM API to
generate trivia questions contextualised to specific POIs along
each trail.
This is a direct implementation of the context-grounded prompting
pattern described in Section~\ref{subsub:sota:llm}: the backend
supplies the model with structured location context and requests
a question appropriate to the POI's historical or cultural
significance.
 
This approach aligns with the current industry best practice for
content-rich, location-aware applications: rather than
pre-authoring questions for each POI (which would be
prohibitively expensive and limit the platform's geographic
scalability), AI generation allows the content layer to scale
with the POI database.
 
The key challenge that the implementation addresses is
\textit{prompt engineering for content reliability}: ensuring that
the model produces questions that are factually grounded,
appropriately difficult, and correctly formatted for the
application's trivia interface~\cite{gravitee2026llmapi}.
The offline limitation of cloud-based LLM calls is also relevant
here: trivia content cannot be generated without network
connectivity, which amplifies the importance of the offline-first
caching capability currently in the backlog.

 
\subsubsection{Summary}\label{subsub:contrast:summary}
 
\begin{table}[ht]
\centering
\small
\renewcommand{\arraystretch}{1.4}
\begin{tabular}{@{} >{\centering\arraybackslash}m{3.2cm}
                    >{\centering\arraybackslash}m{5.8cm}
                    >{\centering\arraybackslash}m{5.8cm} @{}}
\toprule
--
\textbf{Dimension} & \textbf{Alignment} & \textbf{Gap} \\
\midrule
\midrule
--
Location-Based Features &
  Real-time GPS trail interaction; physical presence as core
  interaction primitive &
  Offline-first capability regressed and moved to backlog \\
\midrule
--
Gamification &
  Points and achievement system tied to real POI activities;
  hybrid reward loop &
  No adaptive personalisation; uniform content for all users \\
\midrule
--
LLM Content Generation &
  Context-grounded prompting for POI trivia; scalable content
  without pre-authoring &
  Cloud-only LLM calls require connectivity; no offline
  content fallback \\
\bottomrule
\end{tabular}
\caption{Alignment of internship practices with the state of the art.}
\label{tab:contrast}
\end{table}
 
\noindent
Overall, the internship project demonstrates a strong alignment
with current best practices in location-based mobile application
development, gamified engagement design, and AI-assisted content
generation.
The gaps identified, most notably the offline-first regression, are operational rather than conceptual, and represent natural priorities as the platform moves toward a production release. 